\chapter{练习}

\section{基础}
\begin{homework}[1683 | 查询 第5题]
\begin{lstlisting}[language=SQL]
 select tweet_id
from Tweets
where length(content) > 15;
\end{lstlisting}
数字符串长度是length
\end{homework}

\section{连接}
\begin{homework}[1581 | 连接 第3题]
\begin{lstlisting}[language=SQL]
SELECT customer_id, COUNT(*) AS count_no_trans
FROM Visits
WHERE visit_id NOT IN (
    SELECT visit_id FROM Transactions
)
GROUP BY customer_id
\end{lstlisting}
\begin{dxtips}
    where not in其实就相当于是减法也就是交集，不是减具体数字的都可以用where not in。想不明白先列式子是谁减去谁
\end{dxtips}
\end{homework}

\begin{homework}[197 | 连接 第4题]
\begin{lstlisting}[language=SQL]
SELECT w1.id
FROM Weather w1
JOIN Weather w2
ON DATEDIFF(w1.recordDate, w2.recordDate) = 1
WHERE w1.temperature > w2.temperature;
\end{lstlisting}
\begin{dxtips}[日期题]
    MySQL 中常用的日期函数有两个：
\begin{itemize}
    \item \texttt{DATEDIFF(日期1, 日期2)}：返回日期1 $-$ 日期2的天数差
    \item \texttt{DATE\_ADD(日期, INTERVAL n DAY)}：在某日期上加 $n$ 天
\end{itemize}
\end{dxtips}
\begin{dxtips}
    遇到需要自己表里面数据互相操作最好先连成一张表，如果不需要所有都连接就用子查询生成表的几列当成一个表在链接。弄成两个表就好操作了
\end{dxtips}
\end{homework}

\begin{homework}[1661 | 连接 第5题]
\begin{lstlisting}[language=SQL]
SELECT a1.machine_id, ROUND(AVG(a2.timestamp - a1.timestamp), 3) AS processing_time
FROM Activity AS a1
JOIN Activity AS a2
ON a1.machine_id = a2.machine_id
AND a1.process_id = a2.process_id
AND a1.activity_type <> a2.activity_type
WHERE a2.activity_type = 'end'
GROUP BY a1.machine_id;
\end{lstlisting}
\end{homework}

\begin{definition}[交叉连接（CROSS JOIN）]
交叉连接返回两张表的笛卡尔积，即左表的每一行与右表的每一行两两组合，结果行数为左表行数乘以右表行数。语法为：
\begin{lstlisting}[language=SQL, frame=single, backgroundcolor=\color{white}]
SELECT * FROM 表A CROSS JOIN 表B;
\end{lstlisting}
例如 \texttt{Students}（4行）与 \texttt{Subjects}（3门课）做交叉连接，得到 $4 \times 3 = 12$ 行，即每个学生与每门课的全部组合。常与 \texttt{LEFT JOIN} 配合使用，用于生成"所有可能组合"后再统计实际数据。
\end{definition}

\begin{homework}[1280 | 连接 第7题]
\begin{lstlisting}[language=SQL, frame=single, backgroundcolor=\color{white}]
SELECT s.student_id, s.student_name, sub.subject_name,
       COUNT(e.student_id) AS attended_exams
FROM Students s
CROSS JOIN Subjects sub
LEFT JOIN Examinations e
ON s.student_id = e.student_id
AND sub.subject_name = e.subject_name
GROUP BY s.student_id, s.student_name, sub.subject_name
ORDER BY s.student_id, sub.subject_name;
\end{lstlisting}
\begin{dxtips}
 SELECT 里凡是没有被聚合函数（COUNT、SUM、AVG、MAX、MIN）包裹的列，都必须出现在 GROUP BY 里。因为是要select了我就必须知道他要在那个组里面
\end{dxtips}
\end{homework}

\begin{homework}[570 | 连接 第8题]
\begin{lstlisting}[language=SQL, frame=single, backgroundcolor=\color{white}]
SELECT t.name FROM (
    SELECT e1.name, COUNT(e2.id) AS sum
    FROM Employee AS e1
    JOIN Employee AS e2 ON e1.id = e2.managerId
    GROUP BY e1.id, e1.name  -- {原写法：GROUP BY managerId, e1.name}
) AS t
WHERE t.sum >= 5;
\end{lstlisting}

\begin{dxtips}
    写sql的时候先把最后一步的上一步的表一个长什么样生成对了再进行下一步筛选，如果上一步还复杂就上上步直到出现简单的操作了；
    HAVING 是专门配合 GROUP BY 用的，用来过滤分组后的聚合结果。规则是：
WHERE 过滤的是分组前的原始行
HAVING 过滤的是分组后的聚合值
\end{dxtips}
\begin{lstlisting}[language=SQL, frame=single, backgroundcolor=\color{white}]
SELECT e1.name
FROM Employee AS e1
JOIN Employee AS e2 ON e1.id = e2.managerId
GROUP BY e1.id, e1.name
HAVING COUNT(e2.id) >= 5;
\end{lstlisting}
\end{homework}

\section{函数}
\begin{homework}[1661 | 连接 第5题]
\begin{dxtips}
    这题我的想法是自己需要行与行之间做差所以先自链接，链接需要end对start 然后其他的id需要一致。然后需要做差也就是end-start也就是a2-a1需要a2对type是end最后groupby machine_id求和
\end{dxtips}
\begin{lstlisting}[language=SQL]
SELECT machine_id,
       ROUND(AVG(CASE WHEN activity_type = 'end'
                      THEN timestamp
                      ELSE -timestamp END), 3) AS processing_time
FROM Activity
GROUP BY machine_id;
\end{lstlisting}
\end{homework}

\begin{note}
\texttt{CASE WHEN} 是 SQL 中的条件表达式，相当于编程语言中的 \texttt{if-else}，可用于 \texttt{SELECT}、\texttt{WHERE}、\texttt{ORDER BY} 等子句中。语法为：
\begin{lstlisting}[language=SQL]
CASE WHEN 条件1 THEN 值1
     WHEN 条件2 THEN 值2
     ELSE 值3
END
\end{lstlisting}
例如在聚合函数中，将 \texttt{end} 的时间戳取正、\texttt{start} 的时间戳取负，直接用 \texttt{AVG} 即可得到平均运行时间：
\begin{lstlisting}[language=SQL]
AVG(CASE WHEN activity_type = 'end'
         THEN timestamp
         ELSE -timestamp END)
\end{lstlisting}
\end{note}
\begin{dxtips}
    sql里面的if else语句
\end{dxtips}

\begin{homework}[1934 | 连接 第9题]
\begin{lstlisting}[language=SQL, frame=single, backgroundcolor=\color{white}]
SELECT action,
       CASE WHEN action = 'confirmed' THEN 1 ELSE 0 END AS is_confirmed
FROM Confirmations;
\end{lstlisting}
\begin{dxtips}
    action那行出来就是状态比如说time out或者confirmed。然后这个新的列叫is confirmed里面就全是0和1.然后avg就是求平均数，round（avg，3）就是保留三位小数的平均数。和计算函数一样一般放select ，后面。有case就必须有end
\end{dxtips}

\begin{lstlisting}[language=SQL, frame=single, backgroundcolor=\color{white}]
SELECT s.user_id, IFNULL(stat, 0) AS confirmation_rate
FROM Signups AS s
LEFT JOIN (
    SELECT c.user_id, ROUND(AVG(CASE WHEN action = 'confirmed'
                                     THEN 1 ELSE 0 END), 2) AS stat
    FROM Confirmations AS c
    GROUP BY c.user_id
) AS cs ON s.user_id = cs.user_id;
\end{lstlisting}
\begin{dxtips}
    先写子查询表，在用join把signup表和子查询表连起来，要注意的是如果要连必须给新的表起一个名字比如cs，然后ifnull(stat,0)就是给stat位置是null的自动填写是0
\end{dxtips}
\end{homework}

\section{聚合}
\begin{homework}[1251 | 聚合函数 第2题]
\begin{lstlisting}[language=SQL, frame=single, backgroundcolor=\color{white}]
SELECT ps.product_id,
       IFNULL(ROUND(SUM(ps.price * us.units) / SUM(us.units), 2), 0) AS average_price
FROM UnitsSold AS us
RIGHT JOIN (SELECT * FROM Prices) AS ps
ON us.product_id = ps.product_id
AND us.purchase_date BETWEEN ps.start_date AND ps.end_date
GROUP BY ps.product_id;
\end{lstlisting}

\begin{note}
相较于初始写法，本题做了三处修改：
\begin{itemize}
    \item \texttt{JOIN} 改为 \texttt{RIGHT JOIN}：以 \texttt{Prices} 为右表，保留所有产品，即使没有销售记录也会出现在结果中
    \item 加上 \texttt{IFNULL(..., 0)}：\texttt{RIGHT JOIN} 后没有匹配的产品 \texttt{average\_price} 为 \texttt{NULL}，用 \texttt{IFNULL} 补0
    \item \texttt{GROUP BY us.product\_id} 改为 \texttt{GROUP BY ps.product\_id}：右表没匹配到的行 \texttt{us.product\_id} 为 \texttt{NULL}，必须按 \texttt{ps.product\_id} 分组才能正确输出所有产品
\end{itemize}
\end{note}
\begin{dxtips}
    select内容一变化group by就必须变
\end{dxtips}
\end{homework}

\begin{homework}[1075 | 聚合函数 第3题]
\begin{example}
查询每个项目员工的平均经验年数：
\begin{lstlisting}[language=SQL, frame=single, backgroundcolor=\color{white}]
SELECT p.project_id, ROUND(AVG(e.experience_years), 2) AS average_years
FROM Project AS p
LEFT JOIN Employee AS e ON p.employee_id = e.employee_id
GROUP BY p.project_id;
\end{lstlisting}

\textbf{第一步：LEFT JOIN}，按 \texttt{employee\_id} 连接两表，每行变成"项目id + 员工经验年数"：
\[
\begin{array}{c|c}
\texttt{project\_id} & \texttt{experience\_years} \\
\hline
1 & 3 \\
1 & 2 \\
1 & 1 \\
2 & 3 \\
2 & 2 \\
\end{array}
\]

\textbf{第二步：GROUP BY project\_id}，按项目分组：
\[
\text{project\_id } 1 \to [3,\ 2,\ 1], \quad \text{project\_id } 2 \to [3,\ 2]
\]

\textbf{第三步：AVG}，对每组求平均：
\[
\text{project\_id } 1 \to \frac{3+2+1}{3} = 2.00, \quad \text{project\_id } 2 \to \frac{3+2}{2} = 2.50
\]

\textbf{第四步：ROUND}，保留两位小数输出。
\end{example}
\begin{dxtips}
    group by和计算的可以不是一个
\end{dxtips}
\end{homework}

\begin{homework}[1633 | 聚合函数 第4题]
\begin{note}[主要是主键有问题]
本题的三处易错点：
第一，\texttt{ROUND} 的位数参数写成 \texttt{0} 而非 \texttt{2}，导致结果无小数；
第二，子查询统计总用户数时用了 \texttt{COUNT(DISTINCT user\_name)}，应为 \texttt{COUNT(user\_id)}，\texttt{user\_name} 不是主键，用错列会导致分母计算错误；
第三，\texttt{Register} 表的主键是 \texttt{(contest\_id, user\_id)} 组合，本身不会重复，\texttt{COUNT} 时无需加 \texttt{DISTINCT}。
\end{note}
\begin{dxtips}
    主键不能是null不能重复可以唯一识别

\end{dxtips}
\end{homework}

\begin{homework}[1211 | 聚合函数 第5题]
\begin{note}
MySQL 中，条件表达式（比如 \texttt{rating < 3}）直接放在数值计算里，会自动返回 \texttt{1}（真）或 \texttt{0}（假）。

例如：
\begin{lstlisting}[language=SQL]
SELECT rating < 3 FROM Queries;
\end{lstlisting}
结果就是一列 0 和 1。

所以 \texttt{AVG(rating < 3)} 等价于：
\begin{lstlisting}[language=SQL]
SUM(CASE WHEN rating < 3 THEN 1 ELSE 0 END) / COUNT(*)
\end{lstlisting}

这是 MySQL 特有的行为，标准 SQL 和其他数据库（如 PostgreSQL、SQL Server）不支持，要用 \texttt{CASE WHEN} 才行。力扣默认 MySQL 所以能用，但面试手写 SQL 最好还是写 \texttt{CASE WHEN}，更通用。
\end{note}
\end{homework}

\begin{definition}[DATE\_FORMAT 函数]
\texttt{DATE\_FORMAT(date, format)} 将日期按指定格式转换为字符串，常用格式符如下：
\begin{itemize}
    \item \texttt{\%Y}：四位年份，如 \texttt{2019}
    \item \texttt{\%m}：两位月份，如 \texttt{01}
    \item \texttt{\%d}：两位日期，如 \texttt{07}
\end{itemize}
例如提取年月：
\begin{lstlisting}[language=SQL]
SELECT DATE_FORMAT(trans_date, '%Y-%m') AS month FROM Transactions;
\end{lstlisting}
结果将 \texttt{2019-01-07} 转换为 \texttt{2019-01}，常用于按月分组：
\begin{lstlisting}[language=SQL, backgroundcolor=\color{white}, frame=single]
GROUP BY DATE_FORMAT(trans_date, '%Y-%m')
\end{lstlisting}
\end{definition}

\begin{homework}[1193 | 聚合函数 第6题]
\begin{lstlisting}[language=SQL]
SELECT DATE_FORMAT(trans_date, '%Y-%m') AS month,
       country,
       COUNT(state) AS trans_count,
       SUM(amount) AS trans_total_amount,
       SUM(state = 'approved') AS approved_count,
       SUM(CASE WHEN state = 'approved' THEN amount ELSE 0 END) AS approved_total_amount
FROM Transactions
GROUP BY country, month;
\end{lstlisting}
\begin{dxtips}
    sum(state='approved')自动能给state赋值布尔值是是1，不是赋0，然后就可以sum。然后case 后面要跟着when，then 后面也可以写变量名
\end{dxtips}
\begin{note}
两者本质不同：
\begin{itemize}
    \item \texttt{SUM(state='approved')}：利用 MySQL 把布尔表达式自动转成 0/1，只能用于\textbf{计数}。
    \item \texttt{SUM(CASE WHEN state='approved' THEN amount ELSE 0 END)}：\texttt{THEN} 后面跟的是真实字段值 \texttt{amount}，可以是任意数值，用于\textbf{按条件求和}。
\end{itemize}
所以不能混用——\texttt{SUM(state='approved' * amount)} 这种写法是错的，布尔值和字段值是两回事，必须用 \texttt{CASE WHEN} 才能在条件成立时取字段值。
\end{note}
\end{homework}

\begin{definition}[MIN / MAX 函数]
\texttt{MIN} 和 \texttt{MAX} 是聚合函数，可作用于数值、字符串和日期类型。
\begin{itemize}
    \item \texttt{MIN(col)}：返回该列的最小值，用于日期时返回最早日期
    \item \texttt{MAX(col)}：返回该列的最大值，用于日期时返回最晚日期
\end{itemize}
MySQL 中日期按 \texttt{YYYY-MM-DD} 格式存储，从左到右依次比较年、月、日，因此 \texttt{MIN}、\texttt{MAX}、\texttt{>}、\texttt{<} 对日期均可直接使用。例如：
\begin{lstlisting}[language=SQL]
SELECT customer_id, MIN(order_date) AS first_order_date
FROM Delivery
GROUP BY customer_id;
\end{lstlisting}
返回每位顾客最早的下单日期，即首次订单日期。
\end{definition}

\begin{homework}[1174 | 聚合函数 第7题]
\begin{note}
\texttt{WHERE order\_date IN (子查询)} 中子查询返回的是每个顾客的最早日期，但如果两个不同顾客恰好在同一天下了首次订单，\texttt{IN} 会把两个人的那天所有订单都选进来，不只是首次订单那一行。更安全的写法是用 \texttt{(customer\_id, order\_date)} 组合匹配：
\begin{lstlisting}[language=SQL]
SELECT ROUND(AVG(order_date = customer_pref_delivery_date) * 100, 2)
       AS immediate_percentage
FROM Delivery
WHERE (customer_id, order_date) IN (
    SELECT customer_id, MIN(order_date)
    FROM Delivery
    GROUP BY customer_id
);
\end{lstlisting}
\end{note}

\begin{lstlisting}[language=SQL]
SELECT ROUND(AVG(order_date = customer_pref_delivery_date) * 100, 2)
       AS immediate_percentage
FROM Delivery
WHERE (customer_id, order_date) IN (
    SELECT customer_id, order_date
    FROM Delivery
    GROUP BY customer_id
    HAVING order_date = MIN(order_date)
);
\end{lstlisting}
\begin{dxtips}
    having还是不熟练
\end{dxtips}
\end{homework}

\begin{definition}[DATE\_ADD 与 INTERVAL]
\texttt{DATE\_ADD(date, INTERVAL n unit)} 在日期上加上指定时间间隔，常用单位如下：
\begin{itemize}
    \item \texttt{DAY}：天，如加一天
    \item \texttt{MONTH}：月，如加一个月
    \item \texttt{YEAR}：年，如加一年
\end{itemize}

\end{definition}

\begin{example}
    例如在日期上加一天：
\begin{lstlisting}[language=SQL]
SELECT DATE_ADD('2016-03-01', INTERVAL 1 DAY);
-- 结果：2016-03-02
\end{lstlisting}
也可以用减法：
\begin{lstlisting}[language=SQL]
SELECT DATE_ADD('2016-03-01', INTERVAL -1 DAY);
-- 结果：2016-02-29
\end{lstlisting}
常用于判断某日期的前后一天是否存在记录：
\begin{lstlisting}[language=SQL, backgroundcolor=\color{white}, frame=single]
WHERE event_date = DATE_ADD(first_login, INTERVAL 1 DAY)
\end{lstlisting}
\end{example}

\begin{homework}[550 | 聚合函数 第8题]
\begin{dxtips}
    \begin{itemize}
  \item 函数名写错了，是 \texttt{DATE\_ADD} 不是 \texttt{add\_date}。
  \item \texttt{DATE\_ADD(firstdate, INTERVAL 1 DAY)} 里的 \texttt{firstdate} 是刚起的别名，但 MySQL 不允许在同一层 \texttt{SELECT} 里引用刚定义的别名，需要嵌套一层子查询。
\end{itemize}
\end{dxtips}
\begin{lstlisting}[language=SQL]
SELECT ROUND(
    COUNT(DISTINCT a1.player_id) / (SELECT COUNT(DISTINCT player_id) FROM Activity),
    2
) AS fraction
FROM Activity AS a1
JOIN (
    SELECT player_id, MIN(event_date) AS firstdate
    FROM Activity
    GROUP BY player_id
) AS t ON a1.player_id = t.player_id
AND a1.event_date = DATE_ADD(t.firstdate, INTERVAL 1 DAY);
\end{lstlisting}
\begin{dxtips}
今天的错误是太贪心，想少用子查询。首先应该想明白到底要连接几张表。是这样的，比如最后满足我们要求的只有一个元素，这个需要 \texttt{COUNT} 没问题，但是一旦进行操作以后最开始的表一定变化了，所以一定需要一个最开始表的子查询，也就是 \texttt{(SELECT COUNT(DISTINCT player\_id) FROM Activity)}。除数是一个表，被除数是一个表，求除数再看需要几个表。一个表没法和自己的元素上下判断，所以需要自连接弄成横着的。然后筛选条件一定要写在一起 \texttt{DATE\_ADD(t.firstdate, INTERVAL 1 DAY)}，最后再筛选。
\end{dxtips}
\end{homework}

\begin{note}
\texttt{GROUP BY product\_id} 之后每组变成一行，但 \texttt{sale\_date} 这一组里有多个值，MySQL 不知道该显示哪一个，所以不允许直接 \texttt{SELECT}．\texttt{GROUP BY} 之后 \texttt{SELECT} 的列只能是 \texttt{GROUP BY} 里出现的列或聚合函数（\texttt{MIN}、\texttt{MAX}、\texttt{COUNT}、\texttt{SUM}、\texttt{AVG}）．
\begin{lstlisting}[language=SQL]
SELECT product_id,
       COUNT(*) AS 销售次数,
       MIN(sale_date) AS 最早销售,
       MAX(sale_date) AS 最晚销售
FROM Sales
GROUP BY product_id;
\end{lstlisting}
\end{note}

\begin{note}
\texttt{WHERE} 后面不能带聚合函数，但 \texttt{HAVING} 可以，因为执行顺序不同：
\[
\texttt{FROM} \to \texttt{WHERE} \to \texttt{GROUP BY} \to \texttt{HAVING} \to \texttt{SELECT}
\]
\texttt{WHERE} 在 \texttt{GROUP BY} 之前执行，此时还没有分组，聚合函数根本没算出来，所以不能用．\texttt{HAVING} 在 \texttt{GROUP BY} 之后执行，聚合函数已经算好了，所以可以用．
\end{note}

\begin{homework}[619 | 排序和分组 第3题]
\begin{note}
\textbf{聚合函数的作用域是"组内"，不是"全表"}

一旦有 \lstinline[language=SQL]|GROUP BY|，所有聚合函数（\lstinline[language=SQL]|MAX|、\lstinline[language=SQL]|SUM|、\lstinline[language=SQL]|COUNT| 等）只在\textbf{当前组的行}上计算，不跨组。

\textbf{易错点：} \lstinline[language=SQL]|GROUP BY num| 之后每组已折叠为一行，组内 \lstinline[language=SQL]|num| 全部相同，此时 \lstinline[language=SQL]|MAX(num)| 等于原值，毫无意义。

\textbf{正确做法：}内层分组过滤，外层跨组聚合。

\begin{lstlisting}[language=SQL, backgroundcolor=\color{white}, frame=single]
-- 错误：MAX 只在组内算，等于原值
SELECT MAX(num)
FROM MyNumbers
GROUP BY num
HAVING COUNT(*) < 2;

-- 正确：子查询先得到候选值，外层再取最大
SELECT MAX(num)
FROM (
    SELECT num
    FROM MyNumbers
    GROUP BY num
    HAVING COUNT(*) < 2
) t;
\end{lstlisting}
\end{note}
\end{homework}

\section{高级查询和链接}
\begin{note}
\lstinline[language=SQL]|CASE WHEN| 是 SQL 里的条件判断，相当于 if-else，可放在以下位置：
\begin{itemize}
    \item \lstinline[language=SQL]|SELECT| 后面（最常用，生成新列）
    \item \lstinline[language=SQL]|WHERE| 后面（条件筛选）
    \item \lstinline[language=SQL]|ORDER BY| 后面（按条件排序）
    \item \lstinline[language=SQL]|GROUP BY| 后面（按条件分组）
\end{itemize}
\begin{lstlisting}[language=SQL, backgroundcolor=\color{white}, frame=single]
CASE
    WHEN 条件1 THEN 结果1
    WHEN 条件2 THEN 结果2
    ELSE 默认结果
END
\end{lstlisting}

\end{note}

\begin{homework}[610 | 高级查询和连接 第3题]
\begin{note}
三角形判断条件等价于最大边小于另外两边之和：
\[
\text{GREATEST}(x,y,z) < x+y+z - \text{GREATEST}(x,y,z)
\]
\begin{lstlisting}[language=SQL, backgroundcolor=\color{white}, frame=single]
CASE
    WHEN GREATEST(x, y, z) < x + y + z - GREATEST(x, y, z)
    THEN 'Yes'
    ELSE 'No'
END AS triangle
\end{lstlisting}
\end{note}
\begin{note}
\lstinline[language=SQL]|MAX|/\lstinline[language=SQL]|MIN| 是聚合函数，作用于多行之间的比较（跨行），不能用于同一行内多列的比较。\lstinline[language=SQL]|GREATEST|/\lstinline[language=SQL]|LEAST| 是标量函数，作用于同一行内多个值之间的比较（跨列）。
\end{note}
\end{homework}

\subsection{窗口函数}
\begin{definition}{窗口函数}{}
窗口函数在保留原表所有行的基础上，对每一行计算其"窗口"（周围若干行）内的聚合或排序结果，不像 \lstinline[language=SQL]|GROUP BY| 那样折叠行。基本语法为：
\begin{lstlisting}[language=SQL, backgroundcolor=\color{white}, frame=single]
函数名() OVER (
    PARTITION BY 分组列   -- 可省略
    ORDER BY 排序列
)
\end{lstlisting}
常用窗口函数：
\begin{itemize}
    \item \lstinline[language=SQL]|LAG(col, n)|：当前行往上第 $n$ 行的值
    \item \lstinline[language=SQL]|LEAD(col, n)|：当前行往下第 $n$ 行的值
    \item \lstinline[language=SQL]|RANK()|：当前行的排名（并列跳号）
    \item \lstinline[language=SQL]|ROW_NUMBER()|：当前行的行号（不跳号）
\end{itemize}
\end{definition}
\begin{note}
\texttt{OVER (ORDER BY ...)} 中的排序仅作用于窗口函数的计算过程，相当于一张临时工作视图，不影响最终结果的行顺序。执行流程为：

\begin{lstlisting}[language=SQL, backgroundcolor=\color{white}, frame=single]
原表（无序）
  -> OVER (ORDER BY score DESC)  -- 生成临时排序视图，供窗口函数计算
  -> 窗口函数在临时视图上计算结果
  -> 结果贴回原表行，按原表顺序输出
\end{lstlisting}


\end{note}
\begin{dxtips}
    窗口函数是自连接的高级替代品，覆盖了自连接的常见用途，同时能表达自连接表达不了的"相邻行"语义，性能也更好。
\end{dxtips}
\begin{note}
\texttt{OVER (ORDER BY ...)} 后可跟多列，逗号分隔，规则与普通 \lstinline[language=SQL]|ORDER BY| 相同。常用窗口函数语义如下：

\begin{itemize}
    \item \lstinline[language=SQL]|RANK()|：按排序结果给排名，并列则同名次，下一名跳号
    \item \lstinline[language=SQL]|ROW_NUMBER()|：按排序结果给行号，并列也强制不同，严格连续
    \item \lstinline[language=SQL]|LAG(col, n)|：当前行在排序视图中，往上第 $n$ 行的 \texttt{col} 值；第一行往上无数据则为 \texttt{NULL}
    \item \lstinline[language=SQL]|LEAD(col, n)|：当前行在排序视图中，往下第 $n$ 行的 \texttt{col} 值；最后一行往下无数据则为 \texttt{NULL}
\end{itemize}
\end{note}
\begin{dxtips}
    where在select前面所以不能用select 新造的名字在where里面
\end{dxtips}

\begin{homework}[1164 | 高级查询和连接 第5题]
\begin{lstlisting}[language=SQL, backgroundcolor=\color{white}, frame=single]
select product_id,new_price from Products where( product_id,change_date) in(select product_id,max(change_date)as new
from Products where change_date < '2019-08-16' group by product_id);
--然后先把每一个查询需要提供什么想清楚把子查询就留在上面。然后复制来看有没有对齐什么的
select distinct p1.product_id,ifnull(p2.new_price,10)as price
from Products p1 left join(
select product_id,new_price from Products where( product_id,change_date) in(select product_id,max(change_date)as new
from Products where change_date <= '2019-08-16' group by product_id))as p2
on p1.product_id=p2.product_id
\end{lstlisting}
\begin{dxtips}
IN 只匹配日期值，不匹配产品,所以如果两个都要匹配得弄成数组形式。把做好的子查询留在上面方便查看哪里有问题。
\end{dxtips}
\end{homework}

\begin{note}
普通聚合函数加上 \lstinline[language=SQL]|OVER| 即变为窗口函数，不折叠行，每行保留并附上窗口内的聚合结果。

\begin{lstlisting}[language=SQL, backgroundcolor=\color{white}, frame=single]
SUM()   OVER ()
AVG()   OVER ()
MAX()   OVER ()
MIN()   OVER ()
COUNT() OVER ()
\end{lstlisting}

例如 \lstinline[language=SQL]|SUM(weight) OVER (ORDER BY turn)| 表示按 \lstinline[language=SQL]|turn| 顺序累计求和，每行得到的是"到当前行为止所有人的体重总和"。
\end{note}

\begin{definition}[帧子句 Frame Clause]
帧子句用于在窗口函数中精确定义每一行计算时所参考的行范围，完整语法为：
\begin{lstlisting}[language=SQL]
{ROWS | RANGE | GROUPS} BETWEEN <起点> AND <终点>
\end{lstlisting}
其中起点与终点可取 \texttt{UNBOUNDED PRECEDING}（分区首行）、\texttt{N PRECEDING}（前$N$行）、\texttt{CURRENT ROW}（当前行）、\texttt{N FOLLOWING}（后$N$行）、\texttt{UNBOUNDED FOLLOWING}（分区末行）。
\end{definition}

\begin{homework}[1321 | 高级查询和连接 第11题]
\begin{example}
\texttt{ROWS}：按物理行数计，计算当前行及其前6行共7行的消费总额。
\begin{lstlisting}[language=SQL]
SUM(amount) OVER (
    ORDER BY visited_on
    ROWS BETWEEN 6 PRECEDING AND CURRENT ROW
)
\end{lstlisting}
\end{example}

\begin{example}
\texttt{RANGE}：按值范围计，将所有与当前行 \texttt{visited\_on} 相同的行视为同一帧，计算当天及前6天的消费总额。
\begin{lstlisting}[language=SQL]
SUM(amount) OVER (
    ORDER BY visited_on
    RANGE BETWEEN INTERVAL 6 DAY PRECEDING AND CURRENT ROW
)
\end{lstlisting}
\end{example}

\begin{example}
\texttt{GROUPS}：按组数计，将 \texttt{ORDER BY} 值相同的行视为一组，计算当前组及前2组的消费总额。
\begin{lstlisting}[language=SQL]
SUM(amount) OVER (
    ORDER BY visited_on
    GROUPS BETWEEN 2 PRECEDING AND CURRENT ROW
)
\end{lstlisting}
\end{example}
\end{homework}

\begin{homework}[626 | 高级查询和连接 第9题]
\begin{solution}
按新 $id$ 升序返回每两个连续学生的交换结果。

\begin{lstlisting}[language=SQL]
SELECT
    CASE
        WHEN id % 2 = 1 AND id = (SELECT MAX(id) FROM Seat) THEN id
        WHEN id % 2 = 1 THEN id + 1
        ELSE id - 1
    END AS id,
    student
FROM Seat
ORDER BY id;
\end{lstlisting}

\begin{note}
\texttt{CASE WHEN} 按分支顺序逐一判断，匹配即停止，因此第一个条件必须是最特殊的情形。本题中须先排除"最后一个奇数行"这一特例，再处理一般奇数行；若将两个条件对调，最后一行会错误地执行 $id+1$，导致结果越界。
\end{note}
\end{solution}
\begin{dxtips}
    这个题的关键就是通过调整id来调整名单
\end{dxtips}
\end{homework}

\begin{definition}[PARTITION BY]
\texttt{PARTITION BY} 是窗口函数中的分组子句，语法为：
\begin{lstlisting}[language=SQL]
函数名() OVER (PARTITION BY col1, col2, ... ORDER BY col3)
\end{lstlisting}
与 \texttt{GROUP BY} 的区别在于：\texttt{GROUP BY} 将多行压缩为一行，而 \texttt{PARTITION BY} 在保留所有行的前提下，将数据划分为若干分组，窗口函数在每个分组内部独立计算。不指定 \texttt{PARTITION BY} 时，窗口函数的计算范围为整张表。
\end{definition}
\begin{definition}[DENSE\_RANK()]
\texttt{DENSE\_RANK()} 是窗口函数，用于在分组内对行进行排名，语法为：
\begin{lstlisting}[language=SQL]
DENSE_RANK() OVER (PARTITION BY col1 ORDER BY col2 DESC)
\end{lstlisting}
与 \texttt{RANK()} 的区别在于：并列时不跳过后续排名。例如两行并列第2，\texttt{RANK()} 下一名为第4，\texttt{DENSE\_RANK()} 下一名为第3。
\end{definition}

\subsection{union all}
\begin{definition}[\texttt{UNION ALL}]
将多个 \texttt{SELECT} 语句的结果集纵向拼接，要求各子查询的列数与对应列的数据类型一致。
\begin{lstlisting}[language=SQL]
SELECT col1, col2 FROM table1
UNION ALL
SELECT col1, col2 FROM table2;
\end{lstlisting}
\texttt{UNION ALL} 保留所有行（含重复）；\texttt{UNION} 额外执行去重，性能较差。列名以第一个 \texttt{SELECT} 为准。
\end{definition}
但是为了追求运行效率还是
leftjoin一个你需要所有元素的基础表比较好。

\begin{homework}[1907 | 高级查询和连接 第7题]
\begin{definition}[\texttt{LEFT JOIN} 强制保留分类]
当 \texttt{GROUP BY} 无法保证所有分类出现时，先用 \texttt{UNION ALL} 构造枚举模板表作为左表，再 \texttt{LEFT JOIN} 实际统计结果，缺失分类得到 \texttt{NULL}，用 \texttt{COALESCE} 转为 0。
\begin{lstlisting}[language=SQL]
SELECT category, COALESCE(accounts_count, 0) AS accounts_count
FROM (
    SELECT 'Low Salary'     AS category UNION ALL
    SELECT 'Average Salary' AS category UNION ALL
    SELECT 'High Salary'    AS category
) AS categories
LEFT JOIN (
    SELECT
        CASE
            WHEN income < 20000                 THEN 'Low Salary'
            WHEN income BETWEEN 20000 AND 50000 THEN 'Average Salary'
            ELSE                                     'High Salary'
        END AS category,
        COUNT(*) AS accounts_count
    FROM Accounts
    GROUP BY category
) AS counts USING (category);
\end{lstlisting}
\end{definition}
\end{homework}

想把两个结果搜好了再竖着拼起来最后输出
\begin{lstlisting}[language=SQL]
SELECT name AS results
FROM (
    结果1
) AS t1
UNION ALL
SELECT title AS results
FROM (
    结果2
) AS t2;
\end{lstlisting}
\begin{dxtips}
\begin{itemize}
    \item \texttt{JOIN}：横着拼，增加列
    \item \texttt{UNION ALL}：竖着拼，增加行
\end{itemize}
\end{dxtips}

\begin{homework}[602 | 高级查询和连接 第12题]
    \begin{note}
\texttt{UNION ALL} 的好处是可以通过 \texttt{SELECT} 的时候重新赋予名字（列名），两个子查询如果赋予相同的名字，两个表的数据就能上下排列，方便计算。

\begin{lstlisting}[language=SQL]
SELECT id, SUM(cnt) AS num
FROM (
    SELECT requester_id AS id, COUNT(accepter_id) AS cnt
    FROM RequestAccepted GROUP BY requester_id
    UNION ALL
    SELECT accepter_id AS id, COUNT(requester_id) AS cnt
    FROM RequestAccepted GROUP BY accepter_id
) AS t
GROUP BY id
ORDER BY num DESC
LIMIT 1
\end{lstlisting}
\end{note}
\end{homework}

\section{子查询}
子查询就是创造一个池子以前要写a在1，2，3里面选但是用另一个表表里就是1，2，3。就用select from这个表直接返回1，2，3。
当然作为比较条件也可以是大余小余找第二大的去掉最大的就行。

\begin{homework}[176 | 高级字符串函数 第4题]
\begin{lstlisting}[language=SQL]
SELECT MAX(Salary) AS SecondHighestSalary
FROM Employee
WHERE Salary < (SELECT MAX(Salary) FROM Employee);
\end{lstlisting}
\end{homework}

\begin{definition}[SQL 自定义函数]
自定义函数（User-Defined Function）是用户在数据库中创建的可复用逻辑单元，语法为：

\begin{lstlisting}[language=SQL, frame=single, backgroundcolor=\color{white}]
CREATE FUNCTION 函数名(参数名 参数类型, ...) RETURNS 返回类型
BEGIN
    -- 函数体
    RETURN 返回值;
END
\end{lstlisting}

与存储过程的区别在于：函数必须有返回值，且可以直接嵌套在 \texttt{SELECT} 语句中使用。
\end{definition}

\begin{definition}[SQL 循环]
循环只能写在存储过程或自定义函数内部，MySQL 提供三种循环结构：

\begin{lstlisting}[language=SQL, frame=single, backgroundcolor=\color{white}]
-- 1. LOOP：手动控制退出
loop1: LOOP
    IF 条件 THEN LEAVE loop1; END IF;
END LOOP;

-- 2. WHILE：先判断条件再执行
WHILE 条件 DO
    -- 执行内容
END WHILE;

-- 3. REPEAT：先执行再判断，至少执行一次
REPEAT
    -- 执行内容
UNTIL 条件 END REPEAT;
\end{lstlisting}

\end{definition}

\begin{homework}[177 | 进阶SQL（非SQL50）]
\begin{example}*[177第n高薪水]
    \begin{lstlisting}[language=SQL]
CREATE FUNCTION getNthHighestSalary(N INT) RETURNS INT
BEGIN
    SET N = N - 1;
    RETURN (
        SELECT DISTINCT Salary
        FROM Employee
        ORDER BY Salary DESC
        LIMIT 1 OFFSET N
    );
END
\end{lstlisting}
\end{example}
\end{homework}

\begin{homework}[585 | 高级查询和连接 第13题]
\begin{lstlisting}[language=SQL]
SELECT ROUND(SUM(tiv_2016), 2) AS tiv_2016
FROM Insurance
WHERE tiv_2015 IN (
    SELECT tiv_2015
    FROM Insurance
    GROUP BY tiv_2015
    HAVING COUNT(*) > 1)
AND (lat, lon) IN (
    SELECT lat, lon
    FROM Insurance
    GROUP BY lat, lon
    HAVING COUNT(*) < 2)
\end{lstlisting}
\begin{dxtips}
   只要不是特别简单的筛选，有一个筛选用一个子查询。需要对多列组合 \texttt{GROUP BY} 时直接 \texttt{GROUP BY a, b}；但用 \texttt{WHERE} 筛选多列组合时需要 \texttt{(a, b) IN (SELECT a, b ...)}。
\end{dxtips}
\end{homework}

\section{高级字符串函数 / 正则表达式 / 子句}
\begin{definition}[CONCAT 函数]
\texttt{CONCAT(str1, str2, ...)} 将多个字符串首尾连接，返回拼接结果。

\begin{lstlisting}[language=SQL, frame=single, backgroundcolor=\color{white}]
CONCAT('Hello', ' ', 'World')  -- 结果: 'Hello World'
\end{lstlisting}
\end{definition}

\begin{definition}[LEFT 函数]
\texttt{LEFT(str, n)} 从字符串左边起取 $n$ 个字符。

\begin{lstlisting}[language=SQL, frame=single, backgroundcolor=\color{white}]
LEFT('aLiCe', 1)  -- 结果: 'a'
\end{lstlisting}
\end{definition}

\begin{definition}[SUBSTRING 函数]
\texttt{SUBSTRING(str, pos)} 从第 $pos$ 个字符开始取到末尾（MySQL 下标从 1 开始）。也可写作 \texttt{SUBSTRING(str, pos, len)}，此时只取 $len$ 个字符。

\begin{lstlisting}[language=SQL, frame=single, backgroundcolor=\color{white}]
SUBSTRING('aLiCe', 2)     -- 结果: 'LiCe'
SUBSTRING('aLiCe', 2, 3)  -- 结果: 'LiC'
\end{lstlisting}
\end{definition}

\begin{definition}[UPPER / LOWER 函数]
\texttt{UPPER(str)} 将字符串所有字母转为大写；\texttt{LOWER(str)} 将所有字母转为小写。

\begin{lstlisting}[language=SQL, frame=single, backgroundcolor=\color{white}]
UPPER('aLiCe')  -- 结果: 'ALICE'
LOWER('aLiCe')  -- 结果: 'alice'
\end{lstlisting}
\end{definition}

\begin{homework}[1667 | 高级字符串函数 第1题]
\begin{dxtips}
    我理解一下这种其实就是如果需要对字符串操作需要先确定范围然后拆分操作完了再拼接回去ledt负责顺序
    substring负责确定范围，upper Lower负责操作concat负责拼回去
\end{dxtips}

\begin{dxtips}
    括号多的我们能一行只能一对括号
\end{dxtips}
\end{homework}

\begin{homework}[1527 | 高级字符串函数 第2题]
\begin{lstlisting}[language=SQL]
SELECT patient_id, patient_name, conditions
FROM Patients
WHERE conditions LIKE 'DIAB1%'
   OR conditions LIKE '% DIAB1%';
\end{lstlisting}
\begin{dxtips}
   LIKE 后面字符串要加 \texttt{''}，然后它是标准的检查。

那个发音像的是 \texttt{SOUNDS LIKE}。

还有开头有 \texttt{DIAB1} 就是后面 \texttt{\%}，如果结尾的话前面需要有空格。
\end{dxtips}
\end{homework}

\begin{homework}[196 | 高级字符串函数 第3题]
\begin{dxtips}
    delete就是先select 出要的然后加一个否定的结构
\end{dxtips}
\begin{note}
\begin{lstlisting}[language=SQL]
DELETE p1
FROM Person p1, Person p2
WHERE p1.Email = p2.Email
  AND p1.Id > p2.Id;
\end{lstlisting}

这里的 \texttt{p1} 不是具体表名，而是 \texttt{Person} 的别名；
\texttt{FROM Person p1, Person p2} 表示把同一张表拿出来看两次：
\texttt{Person as p1} 和 \texttt{Person as p2}。

所以 \texttt{DELETE p1} 不是删除一个叫 \texttt{p1} 的表，而是删除
\texttt{Person} 表里那些在 \texttt{p1} 视角下满足条件的行。

可以理解成：\texttt{p1} 是候选被删除的人，\texttt{p2} 是用来比较的人。

当满足 \texttt{p1.Email = p2.Email} 且 \texttt{p1.Id > p2.Id} 时，
说明 \texttt{p1} 这一行邮箱重复，而且 \texttt{Id} 更大，所以删除 \texttt{p1}；
\texttt{p2} 只是参照物，不删除。
\end{note}
\begin{dxtips}
    delete后面只用跟着抽象的表用于确定从那个表删除即可不用具体的列
\end{dxtips}
\begin{lstlisting}[language=SQL]
DELETE FROM Person
WHERE Id NOT IN (
    SELECT Id
    FROM (
        SELECT MIN(Id) AS Id
        FROM Person
        GROUP BY Email
    ) AS t
);
\end{lstlisting}
\begin{note}
\begin{lstlisting}[language=SQL]
SELECT Id FROM (...) AS t
\end{lstlisting}

相当于先把"要保留的 \texttt{Id} 集合"变成一个临时结果 \texttt{t}，再让外层 \texttt{DELETE} 使用它。

简单记：MySQL 里 \texttt{DELETE} 同一张表时，子查询读自己，通常要再套一层临时表。
\end{note}
\end{homework}

\begin{homework}[1484 | 高级字符串函数 第5题]
\begin{definition}[GROUP\_CONCAT]
GROUP\_CONCAT 是 MySQL 中的字符串聚合函数，用来把同一组中的多行字符串合并成一个字符串。

\begin{lstlisting}[language=SQL]
GROUP_CONCAT(DISTINCT product ORDER BY product SEPARATOR ',')
\end{lstlisting}

其中：

\begin{itemize}
  \item DISTINCT product：先对 product 去重；
  \item ORDER BY product：在拼接之前，按照 product 的字典序排序；
  \item SEPARATOR：表示"分隔符"，用来指定拼接时每个字符串之间用什么连接；
  \item SEPARATOR ','：表示用逗号连接不同的 product。
\end{itemize}

所以这句话的意思是：把同一组里的 product 去重后，按字典序排列，再用逗号拼成一个字符串。
\end{definition}
\begin{dxtips}
    group_concat就是能把groupby分组以后组里面的字符串起来，如何第一个product的说明连接什么第二个order是按照什么顺序连接，seprator就是用什么分割
\end{dxtips}
\end{homework}

\begin{homework}[1517 | 高级字符串函数 第7题]
\begin{definition}[REGEXP 正则匹配]
REGEXP 是 SQL 中的正则表达式匹配，用来判断一个字符串是否符合某种模式。

\begin{lstlisting}[language=SQL]
REGEXP '^[A-Za-z][A-Za-z0-9_.-]*@leetcode\\.com$'
\end{lstlisting}

这句话表示筛选出符合下面格式的邮箱：

\begin{itemize}
  \item REGEXP：使用正则表达式进行匹配；
  \item \verb|^|：表示字符串开头；
  \item \verb|[A-Za-z]|：表示第一个字符必须是字母；
  \item \verb|[A-Za-z0-9_.-]*|：表示后面可以有字母、数字、下划线、点和横杠，出现 0 次或多次；
  \item \verb|@leetcode\\.com|：表示必须包含 @leetcode.com，其中点号需要转义；
  \item \verb|$|：表示字符串结尾。
\end{itemize}

所以整体意思是：邮箱必须以字母开头，中间可以包含字母、数字、下划线、点和横杠，并且必须以 @leetcode.com 结尾。
\end{definition}
\end{homework}
